% =============================================================================
%  vsim_bridge_doc.tex
%  Theoretical--Simulation Bridge Project
%  Formation-First VSIM . Identity Entropy Outputs . Scale-Residual Tracking
% =============================================================================
\documentclass[11pt,a4paper]{article}
\usepackage{vsim_common}
\hypersetup{pdftitle={Theoretical--Simulation Bridge Project},pdfsubject={VSIM Internal v5.2.0-internal}}
\begin{document}
\thispagestyle{empty}

\begin{center}
  \eyebrow{VSIM Internal Technical Document}\\[8pt]
  {\LARGE\sffamily\bfseries Theoretical--Simulation Bridge Project}\\[6pt]
  {\normalsize\sffamily
    Formation-First VSIM \quad\textperiodcentered\quad
    Identity Entropy Outputs \quad\textperiodcentered\quad
    Scale-Residual Tracking}
\end{center}

\vspace{10pt}
\noindent
\begin{tabular}{@{}L{2.8cm} L{9cm}@{}}
  \code{Status}   & Internal theoretical planning document \\
  \code{Audience} & VSEPR-SIM / VSIM developers \\
  \code{Purpose}  & Bridge between theoretical identity/scale notation and simulation outputs \\
  \code{Target}   & \texttt{v5.2.0-internal} \\
\end{tabular}

\vspace{10pt}
\begin{warnbox}
\textbf{\small\sffamily\MakeUppercase{Warning}}\quad
This is not a public-facing claims document. It is a design bridge. No need to dress it
up in fake certainty like a grant proposal written by a committee of sleep-deprived raccoons.
\end{warnbox}

\noindent\rule{\linewidth}{2pt}\vspace{6pt}
\tableofcontents
\newpage

% =============================================================================
\section{Project Intent}
% =============================================================================

This chapter formalizes the transition from VSIM as a simulation scripting language into
VSIM as a formation-history and identity-information language. The immediate goal is to
connect theoretical identity notation, scale coalescence, residual structure, and recoverable
information to concrete simulation artifacts. The system should not merely output positions,
velocities, energies, and rendered views. It should also output what the simulation can infer
about how a material formed, what information survived projection across scales, what residual
structure was hidden during packing, and how much identity content remains recoverable after
formation, interaction, or coalescence.

The bridge must remain disciplined: theoretical identity variables belong in derived output
sidecars, overlays, reports, and analysis layers --- not in mutable ground-truth state. VSIM
already separates runtime intent, output records, render state, and validation behavior; render
state is explicitly not supposed to mutate truth state, which gives this project the right
architectural boundary instead of letting theory seep into the kernel like coolant into a bad weld.

% =============================================================================
\section{Core Principle}
% =============================================================================

The core rule is:

\begin{principlebox}
\itshape
Every material simulation must preserve formation context, and every identity-information
output must remain derived from state history rather than written back into state.
\end{principlebox}

Formation describes how the object entered its starting condition. Identity-entropy outputs
describe what is lost, preserved, hidden, or recoverable as the simulation evolves. These
outputs are not replacements for physics. They are an interpretive bridge between the raw
simulation state and the theoretical identity framework.

This follows the existing VSIM direction where parsed sections are transformed into runtime
objects through an intent bridge, including material, environment, and run configuration.
Formation handling already exists as a planned or partially wired runtime concept through
formation stages and field-ramp evaluation, with broader runtime execution deferred toward v5.2.0.

% =============================================================================
\section{Theoretical Basis}
% =============================================================================

The theoretical side of the bridge comes from the identity-matrix framework.
The identity reference defines the universal object in its compact form:

\begin{eqbox}
\eqlabel{Universal Object --- Compact Form}
\[P_i = (I_i,\; H_i,\; \Omega_w)\]
\eqnote{where $I_i$ is the identity matrix, $H_i$ is the deterministic hash matrix,
and $\Omega_w$ is the world seed.}
\end{eqbox}

The extended form adds residual structure and recoverable information:

\begin{eqbox}
\eqlabel{Universal Object --- Extended Form}
\[P_i = (I_i,\; H_i,\; R_i,\; \xi_i^{\star},\; J_i)\]
\eqnote{$R_i$ = residual structure \quad\textperiodcentered\quad
$\xi_i^{\star}$ = true identity-location \quad\textperiodcentered\quad
$J_i$ = recoverable information / distinguishability}
\end{eqbox}

The same reference defines the deterministic path: identity, hash, world seed, and evolution
rules together define the reproducible trajectory.

\begin{eqbox}
\eqlabel{Deterministic Path}
\[\gamma_i(t) = F(I_i,\; H_i,\; \Omega_w,\; t)\]
\end{eqbox}

For the simulation bridge, the most important theoretical object is $J_i$, because it provides
a scalar measure of how much lower-scale structure survives projection into a higher-scale
representation. The notation reference explicitly distinguishes $J_i$ (an information/recovery
measure belonging to an object) from $Q_n$ (the count of active identities at scale $n$, a
simulation inventory count). That distinction matters: one describes information content, the
other describes population.

% =============================================================================
\section{Formation as Required Simulation Context}
% =============================================================================

Formation data becomes mandatory because a material state without formation context is
under-described. A gas mixture, cast bracket, powder bed, welded pipe, crystal surface, or
bead aggregate does not begin as an abstract coordinate list. It has a preparation path. That
preparation path affects density, defects, residual stress, grain state, packing, surface
roughness, thermal state, and transport behavior.

Minimum formation fields should include:

\begin{tomlframe}[VSIM Script --- Formation Block (minimum)]
\begin{lstlisting}[style=tomlstyle]
[formation]
enabled         = true
method          = "declared"
history         = "tracked"
formation_time  = 0.0

[formation.initial_state]
temperature     = 300.0
pressure        = 101325.0
phase           = "solid"
density         = 1.0
defect_state    = "unknown"
grain_state     = "unknown"
\end{lstlisting}
\end{tomlframe}

This block is not decoration. It is the upstream condition for material interpretation.
The simulation can still run without perfect formation knowledge, but the validation layer
should mark missing formation as \code{incomplete}, especially for material, gas, pipe,
surface, powder, and structural simulations.

% =============================================================================
\section{Identity-Entropy Output Layer}
% =============================================================================

The proposed new output layer is a compact derived sidecar with extension \code{.ie}
(Identity Entropy). This file carries identity-information variables per frame and per
identity-bearing object. It links to \code{.dynx}, manifests, reports, and viewer overlays,
but must not mutate particle state.

\begin{tomlframe}[VSIM Script --- outputs.identity\_entropy block]
\begin{lstlisting}[style=tomlstyle]
[outputs.identity_entropy]
enabled          = true
path             = "output/run.ie"
format           = "table"
precision        = 6
attach_to_viewer = true
write_per_frame  = true
write_summary    = true
variables = [
    "J",
    "dataloss",
    "dataentropy",
    "residual",
    "hidden_activity",
    "projection_loss",
    "recoverability",
    "coalescence_loss",
    "lineage_depth",
    "state_quota",
    "entropy_flux",
    "formation_memory"
]
\end{lstlisting}
\end{tomlframe}

This layer belongs beside analysis outputs, symbolic traces, event logs, dashboards, and
manifest records. Existing VSIM output infrastructure already supports many distinct artifact
types --- analysis JSON, event JSON, symbolic traces, reports, dashboards, manifests, and
audit records --- so this is an extension of the output doctrine rather than a random new
species of file goblin.

% =============================================================================
\section{Proposed Variables}
% =============================================================================

The first official variable set should include at least twelve fields.

\vspace{6pt}
\begin{center}\small
\rowcolors{2}{white}{black!3}
\begin{tabular}{@{}L{3.0cm} L{6.5cm} L{3.2cm}@{}}
  \rowcolor{skybg}
  \textbf{Variable} & \textbf{Meaning} & \textbf{Maps to} \\
  \toprule
  \code{J}                 & Recoverable information / distinguishability         & $J_i$ \\
  \code{dataloss}          & Direct identity information loss                     & $1-J_i$ \\
  \code{dataentropy}       & Entropy-like spread/ambiguity of identity channels   & derived \\
  \code{residual}          & Magnitude of packing or projection residual          & $|R_i|$ \\
  \code{hidden\_activity}  & Activity in channels not visible in projected state  & derived \\
  \code{projection\_loss}  & Loss from mapping full identity-space into observable & derived \\
  \code{recoverability}    & Inverse measure of residual/projection loss          & $1-\text{proj\_loss}$ \\
  \code{coalescence\_loss} & Loss from new higher-scale identity formation        & $R_{\text{coal}}^{(n)}$ \\
  \code{lineage\_depth}    & Formation/coalescence layers behind an object        & derived \\
  \code{state\_quota}      & Count of active identities at a scale                & $Q_n$ \\
  \code{entropy\_flux}     & Time derivative of identity entropy                  & $\dot{S}_{\text{id}}$ \\
  \code{formation\_memory} & Fraction of formation history recoverable            & $M_f$ \\
  \bottomrule
\end{tabular}
\end{center}

\vspace{8pt}\noindent A compact row format:

\begin{databox}
{\ttfamily\footnotesize
\textcolor{gray}{frame id J \; dataloss dataentropy residual hidden\_act proj\_loss recov coal\_loss lin state\_q entr\_flux form\_mem flags}\\
0 \; 17 \; 0.9321 \; 0.0679 \; 0.1142 \; 0.0831 \; 0.2200 \; 0.0714 \; 0.9286 \; 0.0000 \; 2 \; 1 \; 0.0042 \; 0.7810 \; 0\\
1 \; 17 \; 0.9260 \; 0.0740 \; 0.1219 \; 0.0887 \; 0.2351 \; 0.0782 \; 0.9218 \; 0.0025 \; 2 \; 1 \; 0.0049 \; 0.7744 \; 0}
\end{databox}

Human-readable enough to debug. Cryptic enough to look like it came out of a basement at
Los Alamos. Balance, finally.

% =============================================================================
\section{Scale Coalescence and Residual Tracking}
% =============================================================================

The identity reference defines scale coalescence as the formation of a new higher-scale
identity when lower-scale identities satisfy a bond/interactability condition and a collective
binding score. The coalescence gate is:

\begin{eqbox}
\eqlabel{Coalescence Gate}
\[G_{\text{coal}}(n)=\mathbf{1}[m\geq3]\;\cdot\;
\mathbf{1}\!\left[\bar{d}^{(n)}\leq\ell_{\text{bond}}^{(n)}\right]\;\cdot\;
\mathbf{1}\!\left[B_n(C_n)\geq\tau_B\right]\]
\eqnote{$m$ = member count \quad\textperiodcentered\quad
$\bar{d}^{(n)}$ = mean inter-member distance \quad\textperiodcentered\quad
$\ell_{\text{bond}}^{(n)}$ = bond length threshold\\[2pt]
$B_n(C_n)$ = collective binding score \quad\textperiodcentered\quad $\tau_B$ = binding threshold}
\end{eqbox}

When the gate fires, a new identity is projected at the next scale:

\begin{eqbox}
\eqlabel{New Higher-Scale Identity}
\[I^{\star}_{(n+1)}=\Pi_n(C_n)\]
\end{eqbox}

The coalescence residual captures internal structure hidden by the new identity formation:

\begin{eqbox}
\eqlabel{Coalescence Residual}
\[R_{\text{coal}}^{(n)}=C_n-A_n\!\left(I^{\star}_{(n+1)}\right)\]
\eqnote{$R_{\text{coal}}^{(n)}$ = structure hidden by new identity formation.
Whatever the higher-scale projection cannot represent is residual.}
\end{eqbox}

For simulation purposes this becomes an output rule: when a formation, clustering, bonding,
packing, or phase event creates a higher-scale object, emit \code{coalescence\_loss} and
\code{residual} fields into the identity-entropy sidecar. This should not assert that the
simulation has ``discovered'' a real new ontology. Yes, science requires this annoying humility.

% =============================================================================
\section{Formation Memory}
% =============================================================================

Formation memory is the bridge between formation-first scripting and identity-entropy output.

\begin{eqbox}
\eqlabel{Formation Memory}
\[M_f=\frac{\text{recoverable formation descriptors}}{\text{total declared formation descriptors}}\]
\end{eqbox}

Formation descriptors may include:

\vspace{4pt}
\begin{center}\small
\begin{tabular}{@{}lll@{}}
temperature history & pressure history    & phase path \\
cooling rate        & deposition order    & packing sequence \\
grain state         & defect seed         & residual stress \\
surface finish      & mixing pathway      & \\
\end{tabular}
\end{center}
\vspace{4pt}

A high $M_f$ value means the current state still carries measurable evidence of its formation
history. A low value means the formation path has been erased, thermalized, randomized,
projected away, or become hidden in residual structure.

This is especially important for:

\vspace{4pt}
\begin{center}\small
\begin{tabular}{@{}lll@{}}
cast parts & welded joints & powder beds \\
gas mixtures & pipe interiors & surface residence simulations \\
crystal defects & irradiated materials & \\
\end{tabular}
\end{center}
\vspace{4pt}

The obvious engineering value is that two materials with identical composition may behave
differently because they formed differently. Humanity did need several thousand years of
metallurgy to rediscover ``history matters,'' but here we are.

% =============================================================================
\section{Script Integration}
% =============================================================================

A minimal VSIM script exercising the full bridge. The older \code{[export]} output system
is kept intact; the \code{[outputs.identity\_entropy]} layer is additive. The doctrine
matters more than the bracket costume.

\begin{tomlframe}[VSIM Script --- identity\_entropy\_surface\_demo \textperiodcentered\ v5.2.0-internal]
\begin{lstlisting}[style=tomlstyle]
[project]
name        = "identity_entropy_surface_demo"
version     = "v5.2.0-internal"
seed_base   = 72001
determinism = true

[formation]
enabled        = true
method         = "deposition"
history        = "tracked"
formation_time = 0.0

[formation.initial_state]
temperature  = 300.0
pressure     = 101325.0
phase        = "solid"
defect_state = "tracked"
surface_state= "roughened"

[material]
formula   = "C"
structure = "graphite"
phase     = "solid"

[system]
mode  = "surface_interaction"
space = "3d"

[[surface]]
type               = "graphite_layer"
roughness          = 0.05
warm_minimum_beads = 100

[particles.projectile]
count    = 128
type     = "generic_bead"
velocity = [0.0, 0.0, -1.5]
height   = 15.0

[run]
mode          = "md"
max_steps     = 8000
dt_fs         = 1.0
temperature_K = 300.0
converge      = true

[observe]
track = ["contact","residence","embedding","diffusion_axis","energy","identity_entropy"]

[export]
write_xyz = true
write_xyzf = true
write_analysis_json = true
write_events_json = true
write_symbolic_trace_json = true
write_report_md = true
write_manifest_json = true

[outputs.identity_entropy]
enabled          = true
path             = "output/run.ie"
format           = "table"
precision        = 6
attach_to_viewer = true
write_per_frame  = true
write_summary    = true
\end{lstlisting}
\end{tomlframe}

% =============================================================================
\section{C++ Implementation Skeleton}
% =============================================================================

Proposed header and translation unit for the identity entropy output layer.
C++17/20, no external dependencies beyond the standard library.
Namespace: \code{vsim::ie}.

\vspace{8pt}
\begin{cppfilelabel}
\textcolor{cpptype}{\ttfamily\small // vsim/identity\_entropy.hpp}
\end{cppfilelabel}
\begin{lstlisting}[style=cppstyle]
// vsim/identity_entropy.hpp -- Formation-First VSIM, Identity Entropy Output Layer
#pragma once
#include <array>
#include <cmath>
#include <cstdint>
#include <fstream>
#include <stdexcept>
#include <string>
#include <vector>

namespace vsim::ie {

// --- Formation Context -------------------------------------------------------

struct FormationContext {
    std::string method;               // "deposition" | "cast" | "declared"
    std::string history;              // "tracked" | "partial" | "unknown"
    double      formation_time = 0.0;
    double      temperature_K  = 300.0;
    double      pressure_Pa    = 101325.0;
    std::string phase;                // "solid" | "liquid" | "gas" | "plasma"
    std::string defect_state;         // "unknown" | "tracked" | "seeded"
    std::string grain_state;          // "unknown" | "columnar" | "equiaxed"
    double      density        = 1.0;

    // M_f : fraction of formation history still recoverable.
    [[nodiscard]] double formation_memory() const noexcept;
};

// --- Identity Entropy Record -------------------------------------------------
//  One row in the .ie sidecar. All fields DERIVED -- never written back to state.

struct IERecord {
    uint64_t frame           = 0;
    uint64_t identity_id     = 0;

    double   J               = 1.0;  // Recoverable information / distinguishability
    double   dataloss        = 0.0;  // 1 - J
    double   dataentropy     = 0.0;  // Spread/ambiguity of identity channels
    double   residual        = 0.0;  // Packing or projection residual magnitude
    double   hidden_activity = 0.0;  // Activity in non-projected channels
    double   projection_loss = 0.0;  // Loss: identity-space -> observable state
    double   recoverability  = 1.0;  // 1 - projection_loss
    double   coalescence_loss= 0.0;  // Loss from higher-scale identity formation
    uint32_t lineage_depth   = 0;    // Formation / coalescence layers behind object
    uint32_t state_quota     = 1;    // Q_n : active identities at this scale
    double   entropy_flux    = 0.0;  // d(identity entropy)/dt
    double   formation_memory= 1.0;  // Recoverable formation fraction  (M_f)
    uint32_t flags           = 0;

    void sync() noexcept {           // Enforce derived invariants before writing.
        dataloss      = 1.0 - J;
        recoverability = 1.0 - projection_loss;
    }
};

// --- Coalescence Gate --------------------------------------------------------
//  G_coal(n) = 1[m >= 3] * 1[d_bar <= l_bond] * 1[B_n >= tau_B]

struct CoalescenceInput {
    uint32_t member_count  = 0;
    double   mean_distance = 0.0;
    double   bond_length   = 0.0;
    double   binding_score = 0.0;
    double   tau_B         = 0.0;
};

[[nodiscard]] inline bool
coalescence_gate(const CoalescenceInput& in) noexcept {
    return (in.member_count  >= 3)
        && (in.mean_distance <= in.bond_length)
        && (in.binding_score >= in.tau_B);
}

// R_coal(n) = || C_n - A_n( I*_(n+1) ) ||
[[nodiscard]] double coalescence_residual(
    const std::vector<double>& cluster_state,
    const std::vector<double>& projected_identity) noexcept;

// --- .ie File Writer ---------------------------------------------------------

class IEWriter {
public:
    explicit IEWriter(const std::string& path, int precision = 6);
    ~IEWriter();
    IEWriter(const IEWriter&)            = delete;
    IEWriter& operator=(const IEWriter&) = delete;

    void write_header();
    void write(IERecord& rec);   // calls sync() internally
    void write_summary();
    void close();

private:
    std::ofstream         stream_;
    int                   precision_;
    uint64_t              rows_written_   = 0;
    std::vector<IERecord> buffer_;
    bool                  header_written_ = false;
    bool                  closed_         = false;
};

} // namespace vsim::ie
\end{lstlisting}

\vspace{8pt}
\begin{cppfilelabel}
\textcolor{cpptype}{\ttfamily\small // vsim/identity\_entropy.cpp}
\end{cppfilelabel}
\begin{lstlisting}[style=cppstyle]
// vsim/identity_entropy.cpp
#include "identity_entropy.hpp"
#include <algorithm>
#include <iomanip>
#include <numeric>

namespace vsim::ie {

double FormationContext::formation_memory() const noexcept {
    using B = bool;
    const std::array<B, 6> known = {{
        !phase.empty()        && phase        != "unknown",
        !defect_state.empty() && defect_state != "unknown",
        !grain_state.empty()  && grain_state  != "unknown",
        temperature_K > 0.0,
        pressure_Pa   > 0.0,
        density       > 0.0,
    }};
    const double n = static_cast<double>(std::count(known.begin(), known.end(), true));
    return n / static_cast<double>(known.size());
}

double coalescence_residual(
    const std::vector<double>& cluster_state,
    const std::vector<double>& projected_identity) noexcept
{
    if (cluster_state.size() != projected_identity.size()) return -1.0;
    double sq = 0.0;
    for (std::size_t i = 0; i < cluster_state.size(); ++i) {
        const double d = cluster_state[i] - projected_identity[i];
        sq += d * d;
    }
    return std::sqrt(sq);
}

IEWriter::IEWriter(const std::string& path, int precision)
    : precision_(precision)
{
    stream_.open(path, std::ios::out | std::ios::trunc);
    if (!stream_.is_open())
        throw std::runtime_error("IEWriter: cannot open " + path);
    stream_ << std::fixed << std::setprecision(precision_);
}

IEWriter::~IEWriter() { if (!closed_) close(); }

void IEWriter::write_header() {
    stream_ <<
        "frame identity_id J dataloss dataentropy residual hidden_activity "
        "projection_loss recoverability coalescence_loss "
        "lineage_depth state_quota entropy_flux formation_memory flags\n";
    header_written_ = true;
}

void IEWriter::write(IERecord& rec) {
    if (!header_written_) write_header();
    rec.sync();
    stream_
        << rec.frame << ' ' << rec.identity_id << ' '
        << rec.J << ' ' << rec.dataloss << ' ' << rec.dataentropy << ' '
        << rec.residual << ' ' << rec.hidden_activity << ' '
        << rec.projection_loss << ' ' << rec.recoverability << ' '
        << rec.coalescence_loss << ' '
        << rec.lineage_depth << ' ' << rec.state_quota << ' '
        << rec.entropy_flux << ' ' << rec.formation_memory << ' '
        << rec.flags << '\n';
    ++rows_written_;
    buffer_.push_back(rec);
}

void IEWriter::write_summary() {
    if (buffer_.empty()) return;
    const double mean_J =
        std::accumulate(buffer_.begin(), buffer_.end(), 0.0,
            [](double a, const IERecord& r){ return a + r.J; })
        / static_cast<double>(buffer_.size());
    const double min_J =
        std::min_element(buffer_.begin(), buffer_.end(),
            [](const IERecord& a, const IERecord& b){ return a.J < b.J; })->J;
    const auto coal =
        std::count_if(buffer_.begin(), buffer_.end(),
            [](const IERecord& r){ return r.coalescence_loss > 0.0; });
    stream_ << "\n# SUMMARY\n"
            << "# rows        " << rows_written_ << '\n'
            << "# mean_J      " << mean_J        << '\n'
            << "# min_J       " << min_J         << '\n'
            << "# coal_events " << coal          << '\n';
}

void IEWriter::close() {
    if (closed_) return;
    write_summary();
    stream_.close();
    closed_ = true;
}

} // namespace vsim::ie
\end{lstlisting}

\vspace{16pt}\noindent\rule{\linewidth}{0.5pt}
{\small\sffamily\textcolor{gray}{VSIM Internal \textperiodcentered\
Theoretical--Simulation Bridge \textperiodcentered\ v5.2.0-internal
\textperiodcentered\ Not for external distribution}}
\end{document}
